% ============================== CHAPTER 3 ============================
\clearpage
\setlecture{Lecture III}\renewcommand{\tabname}{III}
\chapter{Higher Derivatives}\label{ch:higher}

The gradient tells us how a function tilts. To know how it \emph{bends}, whether a flat
point is the top of a hill or the bottom of a valley, we need second derivatives, and in
general derivatives of every order. This chapter shows that mixed partial derivatives can
be taken in any order, builds an algebra of differential operators, and uses it to write
Taylor's formula in several variables: the best polynomial approximation of a function.

\section{Repeated partial derivatives}

Let $f$ be defined on an open set $U\subset\R^2$. Once $D_1f$ and $D_2f$ exist, they are
functions again, and we may differentiate them:
\[ D_1D_1f,\quad D_1D_2f,\quad D_2D_1f,\quad D_2D_2f . \]
Operators are applied from right to left:
\[
D_1D_2f=\pd{}{x}\left(\pd fy\right),\qquad D_2D_1f=\pd{}{y}\left(\pd fx\right).
\]

\begin{example}{mixed partial derivatives}
Let $f(x,y)=\sin(xy)$. Then $\pd fx=y\cos(xy)$ and $\pd fy=x\cos(xy)$.
\begin{itemize}
  \item $D_2D_1f=\pd{}{y}\big(y\cos(xy)\big)=\cos(xy)-xy\sin(xy)$ (product rule in $y$).
  \item $D_1D_2f=\pd{}{x}\big(x\cos(xy)\big)=\cos(xy)-xy\sin(xy)$ (product rule in $x$).
\end{itemize}
The two mixed partials are equal. This is not a coincidence.
\end{example}

\begin{theorem}[label=thm:schwarz]{Equality of mixed partial derivatives}
Let $f$ be a function of two variables defined on an open set $U\subset\R^2$. Assume that
$D_1f$, $D_2f$, $D_1D_2f$ and $D_2D_1f$ exist and are continuous. Then
\[ D_1D_2f=D_2D_1f . \]
\end{theorem}

\begin{intuition}
Look at the four corners of a tiny rectangle and form the alternating sum
$\Delta=f(\text{NE})-f(\text{SE})-f(\text{NW})+f(\text{SW})$. You can read $\Delta$ as
``change in $x$ of the change in $y$'' or as ``change in $y$ of the change in $x$''. It is
the same number either way. Divide by the area $hk$ and shrink the rectangle: one reading
gives $D_2D_1f$, the other gives $D_1D_2f$.
\end{intuition}

\begin{center}
\begin{tikzpicture}[scale=1.1]
  \draw[bpgrid, step=0.5] (-0.7,-0.7) grid (4.7,2.9);
  \fill[plum, opacity=0.1] (0,0) rectangle (4,2.2);
  \draw[plum, line width=1pt] (0,0) rectangle (4,2.2);
  \node[pdot] at (0,0) {}; \node[pdot] at (4,0) {}; \node[pdot] at (0,2.2) {}; \node[dot] at (4,2.2) {};
  \node[font=\small, below left, text=plumdark] at (0,0) {$+f(x,y)$};
  \node[font=\small, below right, text=plumdark] at (4,0) {$-f(x{+}h,y)$};
  \node[font=\small, above left, text=plumdark] at (0,2.2) {$-f(x,y{+}k)$};
  \node[font=\small, above right, text=coraldark] at (4,2.2) {$+f(x{+}h,y{+}k)$};
  \draw[navy, <->, line width=0.6pt] (0,-0.45) -- node[below, font=\scriptsize]{$h$} (4,-0.45);
  \draw[navy, <->, line width=0.6pt] (-0.45,0) -- node[left, font=\scriptsize]{$k$} (-0.45,2.2);
  \node[font=\small, text=navy] at (2,1.1) {$\Delta(h,k)$};
\end{tikzpicture}
\end{center}

\begin{proof}
Fix $(x,y)\in U$ and take $h\neq0$, $k\neq0$ small enough that the closed rectangle with
corners $(x,y)$ and $(x+h,y+k)$ lies in $U$. Let
\[ \Delta=f(x+h,y+k)-f(x+h,y)-f(x,y+k)+f(x,y). \]

\textbf{First reading: vary $x$ first, then $y$.} Put $\varphi(t)=f(t,y+k)-f(t,y)$. Then
$\Delta=\varphi(x+h)-\varphi(x)$. By the mean value theorem there is $s_1$ between $x$ and
$x+h$ with
\[ \Delta=\varphi'(s_1)\,h=\big[D_1f(s_1,y+k)-D_1f(s_1,y)\big]h . \]
Now apply the mean value theorem to $t\mapsto D_1f(s_1,t)$ on the interval from $y$ to
$y+k$: there is $s_2$ between $y$ and $y+k$ with
\[ \Delta=D_2D_1f(s_1,s_2)\,kh . \]

\textbf{Second reading: vary $y$ first, then $x$.} Put $\psi(t)=f(x+h,t)-f(x,t)$. Then
$\Delta=\psi(y+k)-\psi(y)=\psi'(t_2)k=\big[D_2f(x+h,t_2)-D_2f(x,t_2)\big]k$ for some $t_2$
between $y$ and $y+k$, and the mean value theorem in the first variable gives $t_1$ between
$x$ and $x+h$ with
\[ \Delta=D_1D_2f(t_1,t_2)\,hk . \]

\textbf{Compare.} Both expressions equal $\Delta$; cancelling $hk\neq0$,
\[ D_2D_1f(s_1,s_2)=D_1D_2f(t_1,t_2). \]
As $(h,k)\to(0,0)$, both points $(s_1,s_2)$ and $(t_1,t_2)$ tend to $(x,y)$ because they are
trapped in the shrinking rectangle. By continuity of the mixed partials, the left side
tends to $D_2D_1f(x,y)$ and the right side to $D_1D_2f(x,y)$. Hence they are equal.
\end{proof}

\begin{extra}
\textbf{The continuity hypothesis matters.} Let
\[ f(x,y)=\frac{xy(x^2-y^2)}{x^2+y^2}\ \ \big((x,y)\neq O\big),\qquad f(O)=0 . \]
For $y\neq0$, $D_1f(0,y)=\lim_{h\to0}\frac{f(h,y)}{h}=\lim_{h\to0}\frac{y(h^2-y^2)}{h^2+y^2}=-y$,
and also $D_1f(0,0)=0$. So $D_2D_1f(0,0)=\frac{d}{dy}(-y)=-1$. In the same way
$D_2f(x,0)=x$, so $D_1D_2f(0,0)=+1$. The mixed partials exist but differ; they are not
continuous at the origin (Question~\ref{q:schwarz} asks you to check the details).
\end{extra}

\subsection{Higher repeated partial derivatives}

For $f(x,y,z)$ with all partial derivatives under consideration existing and continuous,
the theorem lets us swap any two \emph{adjacent} derivatives, for instance
\[ D_3D_1f=D_1D_3f,\qquad D_1D_3(D_1f)=D_3D_1(D_1f). \]
By swapping adjacent pairs again and again (like sorting cards) we can move all the $D_1$'s
to the left, then all the $D_2$'s, then the $D_3$'s. So if $D_1,D_2,D_3$ occur $m_1,m_2,m_3$
times, every ordering equals
\[ \underbrace{D_1\cdots D_1}_{m_1}\ \underbrace{D_2\cdots D_2}_{m_2}\ \underbrace{D_3\cdots D_3}_{m_3}f . \]

\begin{keyidea}
Under continuity assumptions, a repeated partial derivative depends only on \emph{how many
times} each variable is differentiated, not on the order. The same argument works in any
number of variables.
\end{keyidea}

\begin{example}{order does not matter}
Let $f(x,y,z)=x^2y^3z$. Compute $D_1D_2D_3f$ in two orders.
\begin{itemize}
  \item $D_3f=x^2y^3$, then $D_2D_3f=3x^2y^2$, then $D_1D_2D_3f=6xy^2$.
  \item $D_1f=2xy^3z$, then $D_3D_1f=2xy^3$, then $D_2D_3D_1f=6xy^2$.
\end{itemize}
Same result, as predicted.
\end{example}

\section{Partial differential operators}

\subsection{From derivatives to operators}

In one variable write $D=\frac{d}{dx}$, so $D^mf=f^{(m)}$. Iterating gives
$D^mD^nf=D^{m+n}f$, and with $D^0f=f$ this holds for all $m,n\ge0$.

\begin{definition}{Simple differential operator}
In one variable, $D^m$ is a simple differential operator of order $m$. In $n$ variables, if
$D_1,\dots,D_n$ are the partial derivatives in $x_1,\dots,x_n$, then
\[
D_1^{m_1}\cdots D_n^{m_n}\qquad\text{or}\qquad\left(\pd{}{x_1}\right)^{m_1}\cdots\left(\pd{}{x_n}\right)^{m_n}
\]
(with integers $m_i\ge0$) is a \emph{simple differential operator} of \emph{order}
$m_1+\dots+m_n$.
\end{definition}

\begin{example}{an operator of order 10}
For functions of $(x,y,z)$,
\[
D=\left(\pd{}{x}\right)^3\left(\pd{}{y}\right)^5\left(\pd{}{z}\right)^2=\frac{\partial^{10}}{\partial x^3\partial y^5\partial z^2}
\]
has order $3+5+2=10$. In two variables, $D_1^2D_2^5$ has order $7$.
\end{example}

\subsection{Notation}

All of these appear in the literature and mean the same thing:
\[
\left(\pd{}{x}\right)^2f=\pdd fx=D_1^2f=f_{xx},\qquad
\pd{}{x}\left(\pd fy\right)=\pdm fxy=D_1D_2f .
\]

\begin{trap}
Do not confuse $\left(\pd{}{x}\right)^2f=\pdd fx$ (differentiate twice) with
$\left(\pd fx\right)^2$ (differentiate once, then square). For $f(x,y)=x^2y$:
\[ \pdd fx=2y,\qquad\left(\pd fx\right)^2=(2xy)^2=4x^2y^2 . \]
They are different functions.
\end{trap}

\subsection{General differential operators}

For simple operators $D,D'$ and a constant $c$, define $(D+D')f=Df+D'f$ and $(cD)f=c(Df)$.

\begin{definition}{Differential operator}
A \emph{differential operator} is a sum of terms of the form $c\,D_1^{m_1}\cdots D_n^{m_n}$,
where $c$ is a number and the $m_i$ are integers $\ge0$.
\end{definition}

\begin{example}{applying an operator}
Let $D=3\pd{}{x}+5\left(\pd{}{x}\right)^2-\pi\pd{}{x}\pd{}{y}$ and $f(x,y)=\sin(xy)$.

We need three pieces: $\pd fx=y\cos(xy)$, $\pdd fx=-y^2\sin(xy)$, and (from the first example)
$\pd{}{x}\pd fy=\cos(xy)-xy\sin(xy)$. So
\[
Df=3y\cos(xy)-5y^2\sin(xy)-\pi\big(\cos(xy)-xy\sin(xy)\big).
\]
\end{example}

\textbf{Linearity.} For any differential operator, $D(cf)=cDf$ and $D(f+g)=Df+Dg$, because
ordinary partial differentiation has these properties. Constant-coefficient differential
operators are \emph{linear}.

\subsection{Multiplying operators}

\begin{definition}{Product of operators}
For differential operators $D,D'$, the product $DD'$ is defined by $(DD')f=D(D'f)$.
\end{definition}

\begin{example}{multiplying like polynomials}
Let $D=3\pd{}{x}+2\pd{}{y}$ and $D'=\pd{}{x}+4\pd{}{y}$. Expanding as if they were numbers,
\begin{align*}
DD'&=3\left(\pd{}{x}\right)^2+12\pd{}{x}\pd{}{y}+2\pd{}{y}\pd{}{x}+8\left(\pd{}{y}\right)^2\\
&=3\left(\pd{}{x}\right)^2+14\pd{}{x}\pd{}{y}+8\left(\pd{}{y}\right)^2 .
\end{align*}
The middle terms combine only because $\pd{}{x}\pd{}{y}=\pd{}{y}\pd{}{x}$ (Theorem~\ref{thm:schwarz}).
\end{example}

\begin{keyidea}
Differential operators add and multiply exactly like polynomials:
\[ DD'=D'D,\qquad D(D'+D'')=DD'+DD'' . \]
This rests on the commutation of mixed partial derivatives. It is why we can expand
$\left(h\pd{}{x}+k\pd{}{y}\right)^2$ with the binomial formula below.
\end{keyidea}

\begin{definition}{Standard form, homogeneous operator}
An operator written as a sum of terms $cD_1^{m_1}\cdots D_n^{m_n}$ is in \emph{standard form}.
Each term has \emph{degree} $m_1+\dots+m_n$. If all terms have the same degree $m$, the
operator is \emph{homogeneous of degree $m$}.
\end{definition}

For instance $DD'$ above is homogeneous of degree $2$, while
$3\pd{}{x}+5\left(\pd{}{x}\right)^2$ is not homogeneous.

\subsection{Operators applied to monomials}

\begin{example}{a monomial}
Let $f(x,y)=x^3y^2$. Then
\[ D_1f=3x^2y^2,\quad D_1^2f=2\cdot3\,xy^2,\quad D_1^3f=6y^2,\quad D_1^4f=0, \]
and $D_1^3D_2^2f=D_2^2(6y^2)=12=3!\,2!$.
\end{example}

\begin{proposition}[label=prop:mono]{Derivatives of monomials at the origin}
Let $f(x,y)=x^iy^j$. Then $D_1^iD_2^jf=i!\,j!$, and for $(r,s)\neq(i,j)$,
\[ D_1^rD_2^sf(0,0)=0 . \]
\end{proposition}

\begin{proof}
Differentiating $x^i$ exactly $i$ times gives $i!$, and $y^j$ exactly $j$ times gives $j!$.
If $r>i$, differentiating $x^i$ $r$ times gives $0$. If $r<i$, a positive power $x^{i-r}$
remains, which vanishes at $x=0$. The same argument applies to $y$. So unless $r=i$ and $s=j$,
the value at the origin is $0$.
\end{proof}

\begin{reallife}{the operators of physics}
The most famous differential operator is the \emph{Laplacian}
$\Delta=\left(\pd{}{x}\right)^2+\left(\pd{}{y}\right)^2+\left(\pd{}{z}\right)^2$, homogeneous of degree $2$.
Steady temperatures, electrostatic potentials and incompressible flows all satisfy
$\Delta u=0$. The heat equation $\pd ut=\kappa\Delta u$ and the wave equation
$\left(\pd{}{t}\right)^2u=c^2\Delta u$ are built from the same algebra. For example
$u=x^2-y^2$ and $u=e^x\sin y$ satisfy $\Delta u=0$ in the plane: check it!
\end{reallife}

\section{Taylor's formula}

\subsection{The idea: reduce to one variable}

\begin{intuition}
Taylor's formula in one variable approximates $f(a+h)$ by a polynomial in $h$. To do the same
for $f(P+H)$, walk from $P$ to $P+H$ along the straight segment $P+tH$, $0\le t\le1$. Along
that segment $f$ becomes a one-variable function $g(t)$, and we already know Taylor's formula
for $g$. The only work is to compute the derivatives of $g$ with the chain rule.
\end{intuition}

Recall the one-variable formula: for some $c$ between $a$ and $a+h$,
\[
f(a+h)=f(a)+f'(a)h+\frac{f''(a)}{2!}h^2+\dots+\frac{f^{(r-1)}(a)}{(r-1)!}h^{r-1}+\frac{f^{(r)}(c)}{r!}h^r .
\]
Now let $P=(a,b)$, $H=(h,k)$ and
\[ g(t)=f(P+tH)=f(a+th,\,b+tk), \qquad g(0)=f(P),\quad g(1)=f(P+H). \]

\begin{center}
\begin{tikzpicture}[scale=1]
  \draw[bpgrid, step=0.5] (-0.5,-0.5) grid (5.5,3);
  \node[pdot, label={[font=\small]below left:$P$}] (P) at (0.5,0.3) {};
  \node[dot, label={[font=\small]above right:$P+H$}] (Q) at (4.6,2.4) {};
  \draw[vecc] (P) -- (Q);
  \node[odot, label={[font=\scriptsize, text=navy]above left:$P+tH$}] at ($(P)!0.55!(Q)$) {};
  \node[font=\scriptsize, text=navy, align=left] at (4.3,0.4) {$g(t)=f(P+tH)$\\$t\in[0,1]$};
\end{tikzpicture}
\end{center}

\subsection{Derivatives along a line}

With $x=a+th$, $y=b+tk$ we have $dx/dt=h$, $dy/dt=k$, so the chain rule gives
\[ g'(t)=\pd fx h+\pd fy k, \]
the partials being evaluated at $P+tH$. Differentiate once more, applying the same rule to
each of $\pd fx$ and $\pd fy$:
\begin{align*}
g''(t)&=h\left(\pdd fx h+\pdm fyx k\right)+k\left(\pdm fxy h+\pdd fy k\right)\\
&=h^2\pdd fx+2hk\pdm fxy+k^2\pdd fy=\left(h\pd{}{x}+k\pd{}{y}\right)^2f(P+tH).
\end{align*}
We used the equality of mixed partials to merge the two middle terms.

\begin{definition}{The operator $H\cdot\grad$}
For $H=(h_1,\dots,h_n)$, define the differential operator
\[ H\cdot\grad=h_1D_1+\dots+h_nD_n=h_1\pd{}{x_1}+\dots+h_n\pd{}{x_n}. \]
In two variables $H\cdot\grad=hD_1+kD_2$.
\end{definition}

So $g'(t)=(H\cdot\grad)f(P+tH)$ and $g''(t)=(H\cdot\grad)^2f(P+tH)$.

\begin{trap}
$(H\cdot\grad)^rf(P+tH)$ means: \emph{first} apply the operator $(H\cdot\grad)^r$ to $f$,
\emph{then} evaluate the resulting function at $P+tH$. The vector $H$ is a constant here;
it is not differentiated.
\end{trap}

\begin{theorem}[label=thm:gr]{Higher derivatives along a line}
Let $r\ge1$ and let $f$ have continuous partial derivatives of orders $\le r$ on an open set
$U$. Let $P\in U$, $H$ a vector, and $g(t)=f(P+tH)$. Then
\[ g^{(r)}(t)=\big((H\cdot\grad)^rf\big)(P+tH) \]
for all $t$ such that $P+tH\in U$.
\end{theorem}

\begin{proof}
By induction on $r$. The case $r=1$ is the chain rule (and $r=2$ was computed above).
Assume the formula for $r$, and let $\psi=(H\cdot\grad)^rf$, so that $g^{(r)}(t)=\psi(P+tH)$.
The function $\psi$ has continuous first partial derivatives (it involves derivatives of
order $r$ only, and $f$ has continuous derivatives up to order $r+1$ when we treat this
step). Applying the case $r=1$ to $\psi$,
\[ g^{(r+1)}(t)=\frac{d}{dt}\psi(P+tH)=\big((H\cdot\grad)\psi\big)(P+tH)=\big((H\cdot\grad)^{r+1}f\big)(P+tH). \]
This is the formula for $r+1$.
\end{proof}

\subsection{The formula}

\begin{theorem}[label=thm:taylor]{Taylor's formula}
Let $f$ have continuous partial derivatives up to order $r$ on an open set $U$. Let $P\in U$
and $H$ a vector such that the segment $P+tH$, $0\le t\le1$, lies in $U$. Then there is a
number $\tau$ between $0$ and $1$ such that
\begin{align*}
f(P+H)&=f(P)+\frac{(H\cdot\grad)f(P)}{1!}+\dots+\frac{(H\cdot\grad)^{r-1}f(P)}{(r-1)!}\\
&\quad+\frac{(H\cdot\grad)^rf(P+\tau H)}{r!}.
\end{align*}
\end{theorem}

\begin{proof}
Put $g(t)=f(P+tH)$ and apply the one-variable Taylor formula to $g$ between $0$ and $1$:
\[ g(1)=g(0)+g'(0)+\frac{g''(0)}{2!}+\dots+\frac{g^{(r-1)}(0)}{(r-1)!}+\frac{g^{(r)}(\tau)}{r!} \]
for some $\tau\in(0,1)$. By Theorem~\ref{thm:gr}, $g^{(s)}(0)=(H\cdot\grad)^sf(P)$ and
$g^{(r)}(\tau)=(H\cdot\grad)^rf(P+\tau H)$. Since $g(1)=f(P+H)$ and $g(0)=f(P)$, this is the
claimed formula.
\end{proof}

In $(x,y)$ notation:
\begin{align*}
f(a+h,b+k)&=f(a,b)+\frac{1}{1!}\left(h\pd{}{x}+k\pd{}{y}\right)f(a,b)+\dots\\
&\quad+\frac{1}{(r-1)!}\left(h\pd{}{x}+k\pd{}{y}\right)^{r-1}f(a,b)\\
&\quad+\frac{1}{r!}\left(h\pd{}{x}+k\pd{}{y}\right)^rf(a+\tau h,b+\tau k).
\end{align*}
The powers are expanded exactly as binomials:
\begin{align*}
\left(h\pd{}{x}+k\pd{}{y}\right)^2&=h^2\pdd{}{x}+2hk\frac{\partial^2}{\partial x\partial y}+k^2\pdd{}{y},\\
\left(h\pd{}{x}+k\pd{}{y}\right)^3&=h^3\left(\pd{}{x}\right)^3+3h^2k\left(\pd{}{x}\right)^2\pd{}{y}\\
&\quad+3hk^2\pd{}{x}\left(\pd{}{y}\right)^2+k^3\left(\pd{}{y}\right)^3 .
\end{align*}

\begin{recipe}
\textbf{Second-order Taylor expansion at $(a,b)$} (the version used most in exams):
\[
f(a+h,b+k)\approx f+f_xh+f_yk+\tfrac12\big(f_{xx}h^2+2f_{xy}hk+f_{yy}k^2\big),
\]
all derivatives evaluated at $(a,b)$.
\begin{enumerate}
  \item Compute $f$, the two first and the three second partials.
  \item Evaluate them at $(a,b)$.
  \item Plug into the template. Keep the factor $2$ on $hk$ and the overall $\tfrac12$.
\end{enumerate}
\end{recipe}

\begin{example}{quadratic terms of $\log(1+x+2y)$ at $(2,1)$}
Let $f(x,y)=\log(1+x+2y)$, so $1+x+2y=5$ at $(2,1)$.

\textbf{Step 1: derivatives.} With $u=1+x+2y$:
\[ f_x=\frac1u,\ f_y=\frac2u,\ f_{xx}=-\frac1{u^2},\ f_{xy}=-\frac2{u^2},\ f_{yy}=-\frac4{u^2}. \]
\textbf{Step 2: evaluate at $u=5$.}
\[ f=\log5,\ f_x=\tfrac15,\ f_y=\tfrac25,\ f_{xx}=-\tfrac1{25},\ f_{xy}=-\tfrac2{25},\ f_{yy}=-\tfrac4{25}. \]
\textbf{Step 3: template.}
\[
f(2+h,1+k)=\log5+\left(\tfrac15h+\tfrac25k\right)+\frac1{2!}\left(-\tfrac1{25}h^2-\tfrac4{25}hk-\tfrac4{25}k^2\right)+\cdots
\]
Note $2hk\cdot\left(-\tfrac2{25}\right)=-\tfrac4{25}hk$.
\end{example}

\subsection{Taylor polynomials at the origin}

With $P=O$ and $H=(x,y)$:
\begin{align*}
f(x,y)&=f(0,0)+D_1f(0,0)x+D_2f(0,0)y\\
&\quad+\frac1{2!}\Big[D_1^2f(0,0)x^2+2D_1D_2f(0,0)xy+D_2^2f(0,0)y^2\Big]\\
&\quad+\frac1{3!}\Big[D_1^3f(0,0)x^3+3D_1^2D_2f(0,0)x^2y\\
&\qquad\qquad+3D_1D_2^2f(0,0)xy^2+D_2^3f(0,0)y^3\Big]+R_4 .
\end{align*}
In general $f(x,y)=f(0,0)+G_1(x,y)+\dots+G_{r-1}(x,y)+R_r$ with $G_d$ homogeneous of degree $d$.

\begin{definition}{Polynomial approximation}
The polynomial $f(0,0)+G_1(x,y)+\dots+G_s(x,y)$ is the \emph{polynomial approximation} (or
Taylor polynomial) of $f$ of degree $\le s$.
\end{definition}

\begin{proposition}{Coefficients of a polynomial}
If $G(x,y)=\sum_{i,j}c_{ij}x^iy^j$, then $D_1^rD_2^sG(0,0)=r!\,s!\,c_{rs}$, so
\[ c_{ij}=\frac{D_1^iD_2^jG(0,0)}{i!\,j!}. \]
\end{proposition}

\begin{proof}
Apply $D_1^rD_2^s$ term by term and evaluate at the origin. By
Proposition~\ref{prop:mono} every monomial $x^iy^j$ with $(i,j)\neq(r,s)$ contributes $0$,
and $x^ry^s$ contributes $r!\,s!$. So $D_1^rD_2^sG(0,0)=c_{rs}\,r!\,s!$.
\end{proof}

Consistently, the binomial expansion of the Taylor terms gives
\[
\frac{(xD_1+yD_2)^m}{m!}=\sum_{i=0}^m\frac{x^iy^{m-i}}{i!\,(m-i)!}D_1^iD_2^{m-i},
\]
and so the Taylor polynomial of degree $\le s$ is
\[
\boxed{\,G(x,y)=\sum_{i+j\le s}c_{ij}x^iy^j,\qquad c_{ij}=\frac{D_1^iD_2^jf(0,0)}{i!\,j!}\,}
\]

\begin{keyidea}
The Taylor polynomial of $f$ up to order $s$ is \emph{the} polynomial that has the same
partial derivatives as $f$ at the origin, up to order $s$.
\end{keyidea}

\begin{example}{$\log(1+x+2y)$ at the origin}
At $(0,0)$, $u=1$, so $f=0$, $D_1f=1$, $D_2f=2$, $D_1^2f=-1$, $D_1D_2f=-2$, $D_2^2f=-4$, and
\[ G(x,y)=x+2y-\tfrac12\big(x^2+4xy+4y^2\big). \]
\end{example}

\begin{center}
\begin{tikzpicture}
\begin{axis}[bp3d, width=10cm, height=8.2cm, xlabel=$x$, ylabel=$y$, zlabel=$z$,
   xmin=-1.8, xmax=1.8, ymin=-1.8, ymax=1.8, zmin=-1.6, zmax=1.3, view={-35}{25}]
  \addplot3[surf, shader=flat, draw=cyanbp, fill=cyanbp, fill opacity=0.15, draw opacity=0.6,
     domain=-1.5:1.5, y domain=-1.5:1.5, samples=13] {1-0.5*(x^2+y^2)};
  \addplot3[plumsurf, domain=-1.5:1.5, y domain=-1.5:1.5, samples=22, opacity=0.8, fill opacity=0.8]
     {cos(deg(x))*cos(deg(y))};
  \node[font=\scriptsize, text=cyanbp, anchor=west] at (axis cs:1.5,-1.5,-1.3) {$1-\tfrac12(x^2+y^2)$};
  \node[font=\scriptsize, text=plum, anchor=south] at (axis cs:0,0,1.05) {$\cos x\cos y$};
\end{axis}
\end{tikzpicture}

{\small\color{navy} $f=\cos x\cos y$ (surface) and its degree-2 Taylor polynomial
$1-\tfrac12(x^2+y^2)$ (blue mesh). They agree closely near the origin and separate further out.}
\end{center}

\subsection{Shortcuts with one-variable expansions}

For composite functions it is often faster to substitute into a known one-variable series.

\begin{example}{$\sin(xy)$}
Since $\sin u=u+R_3(u)$ with $\abs{R_3(u)}\le\abs u^3/3!$, put $u=xy$:
\[ \sin(xy)=xy+R_3(xy),\qquad \abs{\sin(xy)-xy}\le\frac{\abs{xy}^3}{3!}. \]
For instance, for $y\neq0$,
$\abs{\dfrac{\sin(xy)-xy}{y}}\le\dfrac{\abs x^3\abs y^2}{3!}\to0$ as $(x,y)\to(0,0)$.
\end{example}

\begin{example}{$\log(1+x+2y)$ the fast way}
From $\log(1+u)=u-\tfrac12u^2+\dots$ with $u=x+2y$:
\begin{align*}\log(1+x+2y)&=x+2y-\tfrac12(x+2y)^2+\dots\\&=x+2y-\tfrac12\big(x^2+4xy+4y^2\big)+\dots\end{align*}
Same polynomial, with no derivatives computed.
\end{example}

\begin{example}{$e^x\cos y$ to order 2}
Multiply $e^x=1+x+\tfrac12x^2+\dots$ by $\cos y=1-\tfrac12y^2+\dots$ and keep degree $\le2$:
\[ e^x\cos y=1+x+\tfrac12x^2-\tfrac12y^2+\dots \]
Check with derivatives: $D_1^2f(0,0)=1$, $D_2^2f(0,0)=-1$, $D_1D_2f(0,0)=0$.
\end{example}

\subsection{Several variables}

For $H=(h_1,\dots,h_n)$ and $H\cdot\grad=h_1D_1+\dots+h_nD_n$, Theorem~\ref{thm:gr} and
Taylor's formula hold word for word. Only the expansion of $(H\cdot\grad)^r$ involves more
terms (use the multinomial formula). For example, in three variables
\[ (H\cdot\grad)^2=h_1^2D_1^2+h_2^2D_2^2+h_3^2D_3^2+2h_1h_2D_1D_2+2h_1h_3D_1D_3+2h_2h_3D_2D_3 . \]

\begin{reallife}{the error of an approximation}
Engineers linearise models every day: a pendulum's $\sin\theta\approx\theta$, a structure's
small deformations, an option price's ``delta-gamma'' approximation in finance
($\Delta V\approx V_SdS+\tfrac12V_{SS}dS^2+V_t\,dt$). Taylor's formula with remainder is what
tells them how big the error of such an approximation can be.
\end{reallife}

\section{Integral expressions}

The mean value remainder contains an unknown $\tau$ that depends on the point. An integral
form avoids this. Start from the fundamental theorem of calculus applied to
$t\mapsto f(tx,ty)$ on $[0,1]$:
\[ f(x,y)=f(0,0)+\int_0^1\frac{d}{dt}f(tx,ty)\,dt . \]
By the chain rule, $\frac{d}{dt}f(tx,ty)=D_1f(tx,ty)\,x+D_2f(tx,ty)\,y$. Hence
\[ f(x,y)=f(0,0)+\int_0^1\big[D_1f(tx,ty)\,x+D_2f(tx,ty)\,y\big]dt . \]

\begin{proposition}{Smooth decomposition}
If $f$ has continuous first partial derivatives on an open disc centred at $O$, then
\[ f(x,y)=f(0,0)+x\,g_1(x,y)+y\,g_2(x,y), \]
with
\[ g_1(x,y)=\int_0^1D_1f(tx,ty)\,dt,\qquad g_2(x,y)=\int_0^1D_2f(tx,ty)\,dt . \]
\end{proposition}

\begin{itemize}
  \item $g_1$ and $g_2$ are as smooth as $D_1f$ and $D_2f$.
  \item One can differentiate under the integral sign with respect to $x$ and $y$.
  \item There is no undetermined $\tau$ depending on $(x,y)$.
\end{itemize}

\begin{example}{a decomposition}
For $f(x,y)=e^{x+y}$: $D_1f=D_2f=e^{x+y}$, so
\[ g_1=g_2=\int_0^1e^{t(x+y)}dt=\frac{e^{x+y}-1}{x+y}\quad(x+y\neq0), \]
and indeed $1+x\,g_1+y\,g_2=1+(e^{x+y}-1)=e^{x+y}$.
\end{example}

\section{Summary}

\begin{keyidea}
\begin{itemize}
  \item With continuous second partials, $D_1D_2f=D_2D_1f$; the order of differentiation never matters.
  \item Differential operators form an algebra that behaves like polynomials.
  \item $(\partial/\partial x)^2f\neq(\partial f/\partial x)^2$.
  \item Along $g(t)=f(P+tH)$: $g^{(r)}(t)=(H\cdot\grad)^rf(P+tH)$.
  \item Taylor: $f(P+H)=\sum_{s<r}\frac{(H\cdot\grad)^sf(P)}{s!}+\frac{(H\cdot\grad)^rf(P+\tau H)}{r!}$.
  \item Taylor coefficients at $O$: $c_{ij}=D_1^iD_2^jf(0,0)/(i!\,j!)$; one-variable series give shortcuts.
\end{itemize}
\end{keyidea}

\section{Exam-style questions}

\begin{exercise}[1]
Let $f(x,y)=x^3y^2+e^{xy}$. Compute all second-order partial derivatives and check that
the mixed ones agree.
\end{exercise}

\begin{exercise}[1]
Show that $u(x,y)=\ln(x^2+y^2)$ satisfies Laplace's equation $u_{xx}+u_{yy}=0$ on
$\R^2\setminus\{O\}$.
\end{exercise}

\begin{exercise}[1]
Write the operator $\left(\pd{}{x}-\pd{}{y}\right)\left(\pd{}{x}+\pd{}{y}\right)$ in standard
form and apply it to $f(x,y)=x^3y+\sin(x+y)$.
\end{exercise}

\begin{exercise}[2]
Find the Taylor polynomial of degree $2$ of $f(x,y)=\sqrt{x^2+y^3}$ at $(1,2)$, and use it
to approximate $\sqrt{1.02^2+1.97^3}$.
\end{exercise}

\begin{exercise}[2]
Find the degree-3 Taylor polynomial of $f(x,y)=e^x\sin y$ at the origin in two ways: with
derivatives, and by multiplying one-variable series. Deduce $D_1^2D_2f(0,0)$ without
differentiating.
\end{exercise}

\begin{exercise}[2]
Without computing any derivative, find $D_1^2D_2^3f(0,0)$ for $f(x,y)=e^x\sin y$.
\end{exercise}

\begin{exercise}[3]\label{q:schwarz}
Let $f(x,y)=\dfrac{xy(x^2-y^2)}{x^2+y^2}$ for $(x,y)\neq O$ and $f(O)=0$.
\begin{enumerate}[label=(\alph*)]
  \item Show that $D_1f(0,y)=-y$ for all $y$ and $D_2f(x,0)=x$ for all $x$.
  \item Deduce $D_2D_1f(0,0)=-1$ and $D_1D_2f(0,0)=1$.
  \item Which hypothesis of Theorem~\ref{thm:schwarz} fails?
\end{enumerate}
\end{exercise}

\begin{exercise}[3]
Using the integral decomposition, write $f(x,y)=\sin(x+y)$ as $x\,g_1+y\,g_2$ with explicit
$g_1,g_2$, and check your answer.
\end{exercise}
